% CONDITIONAL CANDIDATE ONLY — deliberately not included by manuscript.tex.
% Source: Chinese manuscript Tables 3 and 4, registered in
% ../../data/paper_results/legacy_cn_tables.csv.
% Evidence boundary: legacy point estimates; no reconstructable per-seed
% archive, no error bars, and no pooling with v3_leakage_free results.

\subsection{Legacy evidence from the Chinese source manuscript}
\label{subsec:legacy_cn_evidence}

The Chinese source manuscript reported point estimates for four retained-label
ratios before the current fixed-split and leakage-free protocol was established.
These values are retained as legacy evidence because they document the empirical
basis of the original manuscript, but they are not pooled with the revised
experiments. Figure~\ref{fig:legacy_cn_budget} reproduces the two result sets
directly from a table-indexed source-data file. No error bars or significance
claims are shown because the corresponding per-seed prediction archives were not
available locally.

\begin{figure*}[t]
  \centering
  \includegraphics[width=\textwidth]{../../output/figures/legacy_cn/legacy_cn_budget_results_group.pdf}
  \caption{Legacy label-budget trends reported in the Chinese source manuscript.
  \textbf{a}, MAE across four model families. \textbf{b}, RMSE for the
  architecture--physics factorial comparison. Values are point estimates.}
  \label{fig:legacy_cn_budget}
\end{figure*}

In the legacy MAE comparison, PI-MSCL changed from 1.101\% at $r=1.0$ to
1.134\% at $r=0.3$, an increase of 0.033 percentage points. Over the same range,
the reported increases were 0.141, 0.132, and 0.127 percentage points for
XGBoost, LSTM, and CNN-LSTM, respectively. PI-MSCL also had the lowest reported
MAE at each partially supervised ratio: 1.115\%, 1.137\%, and 1.134\% for
$r=0.7$, 0.5, and 0.3. These point estimates are consistent with the original
claim that the combined model was less sensitive to a reduced SOH-label budget;
however, they do not establish the variability or statistical significance of
the differences.

The legacy factorial comparison showed the lowest RMSE for the multi-scale,
physics-informed configuration at all four ratios. Relative to its multi-scale
counterpart without the physical constraint, its reported RMSE was lower by
0.091, 0.086, 0.106, and 0.224 percentage points as $r$ decreased from 1.0 to
0.3. The larger difference at $r=0.3$ is compatible with the proposed role of
soft monotonicity as structural regularization when direct labels are sparse.
Nevertheless, this interpretation remains conditional on the historical
pipeline and must not be presented as confirmation under the revised protocol.

% Claim--evidence map
% Claim: legacy PI-MSCL point estimates vary less with r than legacy baselines.
% Evidence: Chinese source manuscript Table 4; supported as a historical report.
% Claim: legacy multi-scale + PI has the lowest Table-3 RMSE at each r.
% Evidence: Chinese source manuscript Table 3; supported as a historical report.
% Claim: the same effects hold in the revised fixed-split protocol.
% Evidence: incomplete (4/48 runs); must not be stated from this candidate.
